\documentclass[11pt]{article}
\usepackage{microtype,amsmath,amssymb,amsthm,mathtools,booktabs,array}
\usepackage[margin=1in]{geometry}
\usepackage[colorlinks=true,linkcolor=blue!45!black,citecolor=blue!45!black,urlcolor=blue!45!black]{hyperref}
\usepackage{xcolor}
\hypersetup{pdftitle={Gaussian Moment Vanishing with One Base Variable and Constant Quadratic Fibres},pdfsubject={Gaussian moment vanishing with one base variable and constant quadratic fibres}}
\newtheorem{theorem}{Theorem}[section]
\newtheorem{proposition}[theorem]{Proposition}
\newtheorem{lemma}[theorem]{Lemma}
\newtheorem{corollary}[theorem]{Corollary}
\theoremstyle{definition}\newtheorem{definition}[theorem]{Definition}
\theoremstyle{remark}\newtheorem{remark}[theorem]{Remark}
\newcommand{\E}{\mathbb E}
\newcommand{\C}{\mathbb C}
\newcommand{\tr}{\operatorname{tr}}
\newcommand{\wt}{\operatorname{wt}}
\newcommand{\G}{\mathsf G}
\DeclareMathOperator{\Span}{span}
\title{Gaussian Moment Vanishing with One Base Variable\\and Constant Quadratic Fibres}
\author{Roy van Rijn}\date{17 November 2026}
\begin{document}\maketitle

\begin{abstract}
We study complex Gaussian phases with one unrestricted base variable and
any number of quadratic fibre variables, whose quadratic matrix is
constant. Vanishing of every positive pure moment is equivalent to
polynomiality of every marked moment series and to explicit finite
rational identities. The proof allows arbitrary Jordan blocks: purity
forces each nonzero primary block to be cyclic with an affine source,
and a polynomial precision matrix of determinant one cancels every
marked-period pole. We give a sufficient degree bound and a formal
stationary-phase proof of the leading-coefficient obstruction used in
the reduction. Affine normalization also covers correlated real Gaussian
laws with positive-definite covariance. A counterexample with two base
variables explains why the one-base hypothesis matters.
\end{abstract}

\section{The result and its meaning}

Let $u,Y_1,\ldots,Y_s$ be independent standard Gaussian variables.
Expectation is extended complex-linearly to polynomials. We study
\[
 P(u,Y)=Y^TAY+B(u)^TY+g(u),\qquad A=A^T\in\operatorname{Mat}_s(\C),
\]
where $B$ and $g$ have arbitrary polynomial degrees. Thus the fibre
variables $Y$ occur quadratically with constant quadratic part,
while the dependence on the single ordinary Gaussian base $u$ is
unrestricted. The one-dimensional, nonisotropic base is essential. We use standard
normalization here; Corollary~\ref{cor:affinegaussian} treats other
nondegenerate real Gaussian laws.

For a polynomial marker $Q(u,Y)$, independent of $t$, put
\[
 \Phi_Q(t)=\sum_{m\ge0}\E[QP^m]\frac{t^m}{m!}.
\]
These are formal series: convergence of $e^{tP}$ on the original
real integration space is never assumed. Call $P$ \emph{pure} if
$\Phi_1=1$, equivalently $\E[P^m]=0$ for every $m\ge1$. Call it
\emph{safe} if every $\Phi_Q$ is polynomial, equivalently
$\E[QP^m]=0$ for all sufficiently large $m$, for each fixed $Q$.

\begin{theorem}\label{thm:fullfibre}
For the phases above the following conditions are equivalent:
\begin{enumerate}\setlength{\itemsep}{0pt}
\item $P$ is pure;
\item $P$ is safe;
\item the unit period $\Phi_1$ is polynomial;
\item the finite rational identities below hold.
\end{enumerate}
Write
\[
 D=\det(I-2tA),\qquad C=(I-2tA)^{-1}.
\]
The identities are
\begin{gather}
 \deg g\le2,\qquad \deg_u(B^TCB)\le2,\label{eq:degrees}\\
 \alpha u^2+\beta u+\gamma=tg+(t^2/2)B^TCB,
 \qquad K=1-2\alpha,\label{eq:coefficients}\\
 DK=1,\qquad \gamma+\beta^2/(2K)=0.\label{eq:purity}
\end{gather}
No semisimplicity assumption on $A$ is needed. For a pure member,
every nonzero primary block of $A$ is cyclic, and its projection
of $B(u)$ is affine in $u$. All nonlinear source terms are
confined to the nilpotent zero block.

Let $r_0$ be the nilpotency index of that zero block, or zero if
the block is absent, and let $b=\max(1,\deg B)$, with $b=1$ if
$B=0$. Then every original marker obeys the sufficient bound
\begin{equation}\label{eq:cutoff}
 \deg_t\Phi_Q\le\bigl(r_0+(s+1)b\bigr)\deg Q.
\end{equation}
The coefficient is not asserted optimal.
\end{theorem}

The proof has two obstacles. First, quadratic elimination leaves an
arbitrary-degree phase in $u$. The formal leading-coefficient obstruction,
proved in the final section, forces that effective phase to be quadratic.
This gives the rational identities above. Second, rational marked periods
may still have poles. Their cancellation follows from affine-source
rigidity on the nonzero primary spaces and a joint Gaussian precision
matrix of determinant one. Its inverse is polynomial, so every Wick
contraction is polynomial as well. This argument treats defective blocks
of every size.

The algebraic reduction and cancellation occupy the next three sections.
A scalar version of the cancellation identities makes the mechanism
explicit, and a nonlinear example shows how a nilpotent zero fibre can
coexist with a defective nonzero block. The final section supplies the
formal stationary argument used in the first reduction.

\section{A polynomial unit forces the finite purity identities}

We first specify the formal module bridge. For a phase $\varphi(t,x)$
and scalar normalization $h$, the spatial quotient is
\[
 M=\C(t)[x]/(\partial_x+\varphi_x)\C(t)[x],\qquad
 \nabla[q]=[q_t+(\varphi_t+h'/h)q].
\]
The quotient is by the image of the displayed differential operator,
not by an ordinary ideal. The parameter connection is well defined because
the spatial and parameter twisted derivatives commute. The class $[1]$
is called the \emph{seed}; a functional $J$ is \emph{horizontal} when
$J(\nabla[q])=\partial_tJ([q])$. A rational horizontal quotient takes
values in $\C(t)$ with its ordinary derivative.
The full leading-coefficient theorem proved in the last section says:
if the actual degree is $r\ge2$ and
$r v_\infty(h)>v_\infty([x^r]\varphi)$, its differential seed
submodule has no rational horizontal quotient taking $[1]$ to one.
Algebraic $h$ is allowed after a finite cover. The proof retains
all critical multiplicities, collisions and escaping roots, over the
formal Laurent field. Convergence is unnecessary.

\begin{lemma}[Rational seed restriction]\label{lem:seedfunctional}
A horizontal formal functional taking a seed to one restricts on
the differential submodule generated by that seed to a rational
horizontal quotient. If its seed value is a nonzero rational
function $f$, division by $f$ has the same property for the
connection shifted by $-f'/f$.
\end{lemma}
\begin{proof}
For seed value one the successive derivative values are
$1,0,0,\ldots$. Every rational relation is respected by the
given formal functional, so its restriction takes values in
$\C(t)$ and is well defined. For the last assertion differentiate
$J(q)/f$: it is horizontal for $\nabla-f'/f$.
\end{proof}

Quadratic elimination is exact on the top twisted module. For
\[
 -y^TM(t)y/2+L(t,u)^Ty+V(t,u)
\]
with $M$ symmetric and generically invertible, translating
$y=z+M^{-1}L$ gives phase $V+L^TM^{-1}L/2$ and normalization
$h=(\det M)^{-1/2}$. The moving centre gives spatial exact
terms. Differentiating the constant Gaussian class gives
$-\operatorname{tr}(M^{-1}M')/2=h'/h$. Thus this is an
identification of connection and seed, not just of stationary
equations. In the present ordinary Gaussian fibre $M(0)=I$,
so coefficientwise Wick integration at zero is also legitimate.

\begin{lemma}\label{lem:polynomialunit}
If $\Phi_1=f(t)$ is polynomial, then $f=1$ and
\eqref{eq:degrees}--\eqref{eq:purity} hold. Conversely those
identities give $\Phi_1=1$.
\end{lemma}
\begin{proof}
Eliminating $Y$ gives
\[
 \varphi=tg-u^2/2+(t^2/2)B^TCB,\qquad h=D^{-1/2}.
\]
Suppose its actual degree in $u$ is $r\ge3$. The leading
coefficient is $N(t)/D(t)$ with $N$ a nonzero polynomial
divisible by $t$, because that coefficient vanishes at zero.
If $d=\deg D$, then
\[
 r v_\infty(h/f)-v_\infty(N/D)
  =(r/2-1)d+\deg N+r\deg f>0.
\]
The previous lemma transports the actual formal unit functional
to a rational unit quotient for normalization $h/f$. The full
leading-coefficient theorem excludes it. Thus the effective
phase has degree at most two.

For $j\ge3$ its coefficient equation is
$tg_j+(t^2/2)[u^j]B^TCB=0$. Since $C(0)=I$, its first
coefficient gives $g_j=0$, and then $[u^j]B^TCB=0$.
Now Gaussian integration in $u$ is the exact formal identity
\[
 f=(DK)^{-1/2}\exp(R),\qquad R=\gamma+\beta^2/(2K).
\]
Squaring gives $f^2DK=\exp(2R)$. A rational derivative has
zero residues, whereas a rational logarithmic derivative has
the orders of its zeros and poles as residues. Thus $f^2DK$
is constant and $R$ is constant. At zero both normalizations
are one and $R=0$. But $DK$ is polynomial: the denominators
of $C$ divide $D$. Hence the polynomial identity $f^2DK=1$
forces $f=1$, since $f(0)=1$. This gives \eqref{eq:purity}.
Conversely those identities give the unit series one directly.
\end{proof}

In particular, a polynomial unit period in this family is always one.
This implication uses only the seed, independently of the other markers.

\section{Purity forces affine sources on every nonzero block}

Use the nondegenerate symmetric bilinear form underlying the
Gaussian, rather than a Hermitian form. Primary spaces of its
selfadjoint operator $A$ for distinct eigenvalues are orthogonal
and nondegenerate. For completeness, coprime primary polynomials
give their spectral projectors as polynomials in $A$; these
projectors are selfadjoint and orthogonal. Nondegeneracy of each
summand follows from nondegeneracy of the whole space.

Fix $\lambda\ne0$. Let its primary space have dimension $m$
and largest Jordan block size $\ell$. On this space write
$A=\lambda I+N$, $N^\ell=0$, and $l=1-2\lambda t$. Then
\begin{equation}\label{eq:resolvent}
 C_\lambda=\sum_{j=0}^{\ell-1}
             \frac{(2t)^jN^j}{l^{j+1}}.
\end{equation}
The scalar $K$ has a pole there of order at most $\ell$.
But $K=1/D$ has order exactly $m$. Hence $\ell=m$: there is
one Jordan block, with no restriction on its size.

Let $B_\lambda$ be the source projection and put
$I_j(u)=\langle B_\lambda,N^jB_\lambda\rangle$. Each $I_j$
has degree at most two. Indeed for any coefficient of $u^p$,
$p\ge3$, the rational identity $[u^p]B^TCB=0$ has vanishing
principal part at $t_\lambda=(2\lambda)^{-1}$. In
\eqref{eq:resolvent} its highest pole determines
$[u^p]I_{m-1}=0$, since $t_\lambda\ne0$. Descend through the
pole orders to obtain the same assertion for every $I_j$.
Other eigenvalues and the zero block are regular at this point.

Here is the needed elementary property of a selfadjoint cyclic
nilpotent block. Choose a cyclic vector $v$ and put
$e_i=N^{m-i}v$, $1\le i\le m$. Then $Ne_i=e_{i-1}$,
with $e_0=0$. Its Gram entries vanish for $i+j\le m$ and
all anti-diagonal entries equal
$c=\langle v,N^{m-1}v\rangle$. The basis is nondegenerate,
so its anti-triangular determinant forces $c\ne0$.
Write $B_\lambda=\sum_i z_i(u)e_i$. Then
\begin{align}
 I_{m-1}&=c z_m^2,\\
 I_{m-k}&=2c z_m z_{m+1-k}
              +R_k(z_{m+2-k},\ldots,z_m),\quad 2\le k\le m,
 \label{eq:triangular}
\end{align}
where $R_k$ is a constant quadratic form. These identities
follow by shifting the second basis index by $m-k$: all pairs
with index sum less than $2m+1-k$ vanish, and the two boundary
pairs have coefficient $c$.

Since $I_{m-1}$ has degree at most two, $z_m$ is affine.
It is nonconstant. In fact write $D=l^mD_\lambda$ and
$z_m=b_m u+c_m$. Multiplying $K=1/D$ by $l^m$ and evaluating
at $t_\lambda$ gives the exact top-pole certificate
\begin{equation}\label{eq:toppole}
 -t_\lambda^2(2t_\lambda)^{m-1}
          c b_m^2 D_\lambda(t_\lambda)=1.
\end{equation}
In particular $b_m\ne0$. Induction in \eqref{eq:triangular}
now makes every $z_i$ affine: after the higher coordinates
are affine, $z_mz_i$ has degree at most two, and the nonzero
linear leading coefficient of $z_m$ forces $\deg z_i\le1$.
This proves affine-source rigidity on every nonzero block,
including every defective block.

\section{A polynomial precision matrix cancels every marker pole}

Split $Y=(Y_0,Y_*)$ into the nilpotent zero primary space and
the sum of the nonzero spaces, of dimensions $s-k$ and $k$.
Choose orthonormal Gaussian coordinates within both summands;
they exist over $\C$ since the restricted bilinear forms are
nondegenerate. The preceding Jordan-chain calculation was only
a proof of an invariant affine-source statement, not an
unjustified reset of Gaussian covariances. In these coordinates
$A_0$ is nilpotent, $A_*$ is symmetric and invertible, and
\[
 B_*(u)=b_*u+c_*,\qquad
 C_0=(I-2tA_0)^{-1}\in\C[t],\qquad C_*=(I-2tA_*)^{-1}.
\]
Because $B_*^TC_*B_*$ is quadratic in $u$,
\eqref{eq:degrees} makes $B_0^TC_0B_0$ quadratic as well.
Write the following polynomial in $u,t$ as
\[
 tg+(t^2/2)B_0^TC_0B_0=a(t)u^2+d(t)u+e(t).
\]
Integrating $Y_0$ leaves an ordinary joint Gaussian in
$X=(u,Y_*)$, with precision matrix and linear source
\begin{equation}\label{eq:precision}
 M(t)=\begin{pmatrix}
  1-2a(t)&-t b_*^T\\-t b_*&I-2tA_*
 \end{pmatrix},\qquad
 L(t)=\binom{d(t)}{t c_*}.
\end{equation}
Both $M$ and $L$ are polynomial. The Schur determinant identity and
\eqref{eq:purity} give the central cancellation certificate
\begin{equation}\label{eq:unimodular}
 \det M=D\bigl(1-2a-t^2b_*^TC_*b_*\bigr)=DK=1,
 \qquad M^{-1}=\operatorname{adj}M\in\operatorname{Mat}_{k+1}(\C[t]).
\end{equation}
The equalities include $k=0$, with the empty determinant equal to one.

The normalized tilted $X$ functional therefore has polynomial
mean $M^{-1}L$ and polynomial covariance $M^{-1}$. Its scalar
normalization is exactly one because $\Phi_1=1$. Conditional
on $u$, the eliminated variables $Y_0$ have mean $tC_0B_0(u)$
and covariance $C_0$, both polynomial in $t,u$. Conditional
Wick expansion of any original $Q(u,Y_0,Y_*)$ consequently
gives a polynomial in $t,u,Y_*$. Joint Wick expansion using
\eqref{eq:unimodular} makes its period polynomial in $t$.
This proves purity implies safety for every matrix $A$.
The other implications and the finite classification follow
from Lemma~\ref{lem:polynomialunit} and the definition of safety.

For the degree bound, if the zero block is
present, $\deg C_0\le r_0-1$ and $k+r_0\le s$; when it is
absent take $r_0=0$. All coefficients of $a,d,e$ have degree
at most $r_0+1$. In each cofactor of $M$ the high-degree base
entry can occur at most once. Thus
\[
 \deg M^{-1}\le k+r_0\le s.
\]
Cofactors in the base column do not contain that entry and
have degree at most $k$. Multiplying this column by $d$, and
all other columns by the corresponding entries of $tc_*$,
gives $\deg(M^{-1}L)\le s+1$.

For a marker of total degree $q$, conditional zero-block Wick
expansion has coefficient degree at most $r_0q$ in $t$ and
remaining variable degree at most $bq$. A joint Gaussian
monomial of degree $j$ then has expectation degree at most
$(s+1)j$, since its mean degrees are at most $s+1$ and its
covariance degrees at most $s$. This proves
\eqref{eq:cutoff} for every fixed polynomial marker.

\subsection{Scalar pole-cancellation certificates}

For simple eigenvalues the matrix identity \eqref{eq:unimodular} has a
useful scalar form. It shows directly which numerators cancel the
apparent poles left by conditional Gaussian integration.

\begin{corollary}\label{cor:simplecertificates}
Suppose every nonzero eigenvalue of $A$ is semisimple and $P$ is pure.
Each such eigenvalue is then simple. Write
\[
 B_\lambda(u)=b_\lambda u+c_\lambda,\qquad
 l_\lambda=1-2\lambda t,\quad
 t_\lambda=(2\lambda)^{-1},\quad D_\lambda=D/l_\lambda,
 \qquad \mu=D\beta.
\]
Then $\mu$ is polynomial, $b_\lambda\ne0$, and
\begin{equation}\label{eq:scalarcertificate}
 t_\lambda^2b_\lambda^2D_\lambda(t_\lambda)=-1,
 \qquad \mu(t_\lambda)=-c_\lambda/b_\lambda.
\end{equation}
In particular the tilted mean and variance
\begin{equation}\label{eq:scalarcancellation}
 \E_tY_\lambda=
    \frac{t(b_\lambda\mu+c_\lambda)}{l_\lambda},
 \qquad
 \operatorname{Var}_t(Y_\lambda)=
    \frac{1+t^2b_\lambda^2D_\lambda}{l_\lambda}
\end{equation}
are polynomials. Here $\E_t$ denotes the formal functional
$\E_tQ=\Phi_Q$, normalized by $\Phi_1=1$.
\end{corollary}
\begin{proof}
A semisimple cyclic primary space has dimension one. The only pole
of its resolvent is $1/l_\lambda$, so evaluating $DK=1$ at
$t_\lambda$ gives the first identity in \eqref{eq:scalarcertificate}.
Likewise the polar contribution to $\beta$ there is
$t^2b_\lambda c_\lambda/l_\lambda$. Since every denominator
in $\beta$ divides $D$, the function $\mu=D\beta$ is polynomial,
and evaluation gives its stated value.

After eliminating all fibres, $u$ has tilted mean $\mu$ and variance
$D$. Conditional on $u$, the $\lambda$ coordinate has mean
$t(b_\lambda u+c_\lambda)/l_\lambda$ and variance $1/l_\lambda$.
These give \eqref{eq:scalarcancellation}; both numerators vanish at
$t_\lambda$ by \eqref{eq:scalarcertificate}. The other joint
covariances are polynomial directly, since
\[
 \operatorname{Cov}_t(u,Y_\lambda)=\frac{tDb_\lambda}{l_\lambda},
 \qquad
 \operatorname{Cov}_t(Y_\lambda,Y_\nu)=
 \frac{t^2b_\lambda b_\nu D}{l_\lambda l_\nu}\quad(\lambda\ne\nu).
\]
Thus the same cancellation can be checked coordinate by coordinate.
\end{proof}

\subsection{Correlated Gaussian laws}

\begin{corollary}\label{cor:affinegaussian}
Let $(u,Y)$ have a real Gaussian law with positive-definite joint
covariance, allowing arbitrary means and correlations. For
$P=Y^TAY+B(u)^TY+g(u)$ with constant symmetric $A$, purity, safety
and polynomiality of $\Phi_1$ are equivalent. The finite identities
\eqref{eq:degrees}--\eqref{eq:purity} apply after affine normalization.
With $b=\max(1,\deg B)$, taking $b=1$ if $B=0$, every pure member satisfies
\[
 \deg_t\Phi_Q\le\bigl(s+(s+1)b\bigr)\deg Q.
\]
\end{corollary}
\begin{proof}
Write $\bar u=\E[u]$ and $\nu=\E[Y]$. A block lower-triangular
Cholesky factor of the covariance gives
\[
 u=\bar u+\sigma U,\qquad Y=\nu+vU+LZ,
\]
where $U,Z_1,\ldots,Z_s$ are independent standard Gaussians,
$\sigma>0$, and $L$ is invertible. Substitution preserves the fibre family:
\[
 A'=L^TAL,\qquad
 B'(U)=L^T\bigl(B(\bar u+\sigma U)+2A(\nu+vU)\bigr),
\]
with a polynomial residual $g'(U)$. Thus $b'=\max(1,\deg B')\le b$,
with the same convention if $B'=0$. This invertible affine substitution
preserves total marker degrees and every moment identity. Apply
Theorem~\ref{thm:fullfibre} to the normalized coefficients; its zero-block
index $r'_0$ satisfies $r'_0\le s$, giving the stated bound.
\end{proof}

\section{Which fibre matrices occur?}

\begin{corollary}\label{cor:matrixexistence}
For a fixed complex symmetric matrix $A$, a pure member of this
fibre family exists if and only if every nonzero primary space
of $A$ is cyclic. Whenever it exists, a quadratic member already
exists. Zero primary blocks are unrestricted.
\end{corollary}
\begin{proof}
Necessity was proved above. For sufficiency set the zero-block
source to zero. If there are no nonzero eigenvalues,
$P=Y^TAY$ suffices. Otherwise let $k$ be the total nonzero
multiplicity, $D=\det(I-2tA_*)$, and $a=-\operatorname{tr}A_*$.
The rational function
\[
 H(t)=\frac{1-2at-1/D}{t^2}
\]
is regular at zero since $D^{-1}=1+2t\operatorname{tr}A_*+O(t^2)$.
It is proper at infinity and has poles only at nonzero-eigenvalue
points, of their exact multiplicities. The functions
$(2t)^j/(1-2\lambda t)^{j+1}$, for $0\le j<m_\lambda$,
form a basis of all such proper rational functions: their
principal parts are triangular with nonzero top coefficients.
Expand $H$ uniquely in this basis, with constants $h_{\lambda,j}$.
For each block its last constant $h_{\lambda,m_\lambda-1}$
is nonzero because $H$ has its full pole there.

Choose a constant vector $b_\lambda$ with
$\langle b_\lambda,N^jb_\lambda\rangle=h_{\lambda,j}$.
This is constructive in every block size: the last coordinate
solves $c z_m^2=h_{\lambda,m-1}\ne0$, and then
\eqref{eq:triangular} solves the earlier coordinates successively
by division by $2cz_m\ne0$. Combine these vectors to $b_*$.
The polynomial
\[
 P=Y^TAY+u b_*^TY_*+a u^2
\]
has $B^TCB=u^2H(t)$ and
$D(1-2at-t^2H)=1$. Its remaining scalar identities are
$\beta=\gamma=0$, so it is pure, including for defective blocks
of every allowed size.
\end{proof}

If the zero primary space is absent, every pure member is
itself quadratic, so Liu's sharper quadratic mixed cutoff
applies~\cite[Theorems~2.1 and~2.2]{Liu}. That established
quadratic theorem and its optimal coefficient are not new here.
The all-degree extension concerns unrestricted nonlinear sources
in the nilpotent zero fibre and the uniform pole-cancellation
argument for the remaining blocks.

Nonlinear sources can coexist with defective nonzero blocks. Let $u$ be ordinary Gaussian and let $(Z,W)$ and
$(Z_0,W_0)$ be independent circular Gaussian pairs, with
$\E[ZW]=\E[Z_0W_0]=1$. For every $d\ge2$ set
\[
 P_d=2ZW+W^2+iu(2Z+3W)-2u^2+W_0u^d.
\]
The first fibre block has Gram matrix
$\left(\begin{smallmatrix}0&1\\1&0\end{smallmatrix}\right)$
and operator $I+N$, where
$N=\left(\begin{smallmatrix}0&1\\0&0\end{smallmatrix}\right)$.
Its source vector is $u(3i,2i)$, so
\[
 D=(1-2t)^2,\qquad
 B_*^TC_*B_*=u^2\frac{-12+16t}{(1-2t)^2},\qquad
 DK=(1-2t)^2(1+4t)-t^2(-12+16t)=1.
\]
The zero fibre contributes $W_0u^d$ and has zero source norm.
The remaining scalar exponent is zero. Thus every $P_d$ is
pure and safe, has degree $d+1$, and combines an unrestricted
nonlinear source with a nonzero Jordan block of size two.

\section{Why one base variable matters}

The hypothesis cannot be extended uniformly to two ordinary
Gaussian base variables, even with a semisimple constant fibre
matrix. Long's cubic four-variable counterexample provides the
obstruction~\cite[Proposition~4.1]{Long}. In two independent
circular pairs put
\[
 U=(Z_1,W_1),\quad Y=(Z_2,W_2),\quad h(U)=W_1(1-Z_1).
\]
Then
\[
 P_4=(1+Z_2)(h(U)+W_2)
     =Z_2W_2+h(U)Z_2+W_2+h(U).
\]
It has a constant quadratic fibre part $Z_2W_2$, an arbitrary
polynomial source in the two-dimensional Gaussian base, and
no degenerate Gaussian covariance. Its fibre operator is the
semisimple scalar $I/2$.

Here is the all-order witness. Circular contraction gives
$\E_{U}[e^{s h(U)}]=(1+s)^{-1}$ and, for every polynomial $A$,
\begin{align*}
 \E[A(Z_2)e^{tP_4}]
 &=\frac{A(t/(1-t))}{1-t}\,
       \frac1{1+t(1+t/(1-t))}\\
 &=A(t/(1-t)).
\end{align*}
The contraction follows directly from
$\E[Z^aW^b]=\delta_{ab}a!$ by formal coefficient expansion.
Consequently $\Phi_1=1$ but $\Phi_{Z_2}=t/(1-t)$, and
$\E[Z_2P_4^m]=m!$ for every $m\ge1$. This example and its
failure are Long's result. Its presentation as a constant
quadratic fibre explains the precise limit of the new theorem:
one ordinary base variable is a substantive rigidity hypothesis.

More generally let $\mathsf H(r,s)$ assert that purity always
implies safety for
\[
 P(U,Y)=Y^TAY+B(U)^TY+g(U),\quad
 U\in\C^r,\quad Y\in\C^s,
\]
in independent standard Gaussian coordinates, with $A$ constant
and all polynomial degrees unrestricted. Combining the preceding
theorem with the credited binary, quadratic and counterexample
results gives the following ranges:
\begin{center}
\begin{tabular}{@{}lccc@{}}
\toprule
Base variables & $s=0$ & $s=1$ & $s\ge2$\\\midrule
$r=0,1$ & holds & holds & holds\\
$r=2$ & holds & not resolved here & failures exist\\
$r\ge3$ & failures exist & failures exist & failures exist\\
\bottomrule
\end{tabular}
\end{center}
Here ``holds'' means that every pure phase in the indicated family is safe.
For $r=0$ the phase is quadratic, covered by Liu (or a constant
if $s=0$). For $r=1$ use Theorem~\ref{thm:fullfibre}, adding
an unused fibre variable when $s=0$. For $(r,s)=(2,0)$ use Wilson's
two-variable theorem~\cite[Theorem~1]{Wilson}. Long's example
above covers $(2,2)$. For $(3,0)$ use his quartic
\[
 (1+Z)\bigl(W-(2+Z)T^2/2\bigr),\qquad Q=Z,
\]
which has unit period one and marked period $t/(1-t)$
\cite[Theorem~5.1]{Long}. Ignoring added independent coordinates
extends both failures to all indicated cells.

The unresolved cell is $(r,s)=(2,1)$: does a single constant quadratic
fibre over two ordinary Gaussian base variables preserve safety of pure
phases? The companion paper \emph{Gaussian Marker Reduction with a Constant
Quadratic Fibre} reduces this question to markers in the two base
variables. Those periods have a parameter-dependent binary phase, to
which Wilson's fixed-phase theorem does not directly apply. The table
records the deductions above and makes no claim of exhaustive literature
coverage. Safety implying purity for arbitrary Gaussian phases also
remains outside the present theorem.

\section{Formal stationary phase}
\subsection{A formal leading-coefficient criterion}\label{sec:stationaryballs}

For a polynomial phase, the remaining shrinking argument can be
replaced by a valuation calculation. The construction below includes
critical points that collide, remain multiple, or escape to infinity.
The key is to group them before making a Gaussian approximation.

Write $v$ for the valuation of a formal parameter $r$, normalized by
$v(r)=1$. The valuation of a polynomial is the least valuation of its
coefficients. We use the formal Newton--Puiseux theorem: finitely many
algebraic elements over $\C((r))$ belong to one field
$\C((r^{1/N}))$~\cite[Corollary~1.10]{MannaaCoquand}.
Characteristic zero makes each minimal polynomial separable;
no convergence hypothesis is needed.

\begin{lemma}[Polynomial separation of formal balls]\label{lem:ballseparation}
Let finitely many disjoint balls have centers $c_B$ in a Puiseux field
and radii specified by $v(x-c_B)\geq\rho_B$. After one common
ramification, put $\delta_B=r^{\rho_B}$. For any one ball $B$ and
integer $j\geq0$, there is a polynomial $q_{B,j}(x)$ over the Laurent
field whose substitution $x=c_C+\delta_C y$ has integral coefficients
for every ball $C$, and whose reduction is $y^j$ on $B$ and zero on
every other ball.
\end{lemma}
\begin{proof}
The finitely many valuations $v(c_B-c_C)$ give a finite tree of
clusters of centers. At a node with scale $r^d$, the coordinate
$u=(x-c)/r^d$ has distinct constant reductions $a_1,\ldots,a_s$
on its children. The Lagrange polynomial
\[
 e_i(u)=\prod_{k\ne i}\frac{u-a_k}{a_i-a_k}
\]
reduces to one on child $i$ and zero on every other child.

Start at the chosen leaf with $((x-c_B)/\delta_B)^j$. Induct upward.
Suppose a polynomial $h$ already has the required reductions on the
chosen child. On each of the other child balls its coefficients have
some finite lower valuation. The restriction of $e_i$ there has
strictly positive valuation. Consequently $e_i^M h$, for a sufficiently
large fixed integer $M$, has positive coefficient valuation on all
those children. On the chosen child $e_i^M=1+$ positive-order terms,
so the existing reductions are unchanged. There are finitely many
nodes, and this constructs the required polynomial. A node with just
one leaf requires no filter. This also explains why independently
normalized Lagrange polynomials, with no control of their growth on
other scales, would not suffice.
\end{proof}

\begin{theorem}[Formal stationary balls]\label{thm:stationaryballs}
Let $\varphi(r,x)\in\C((r))[x]$ have degree $D\geq2$, and let
$h(r)\in\C((r))\setminus\{0\}$. Put
\[
 M=\C((r))[x]/(\partial_x+\varphi_x)\C((r))[x],\qquad
 \nabla[q]=[q_r+(\varphi_r+h'/h)q].
\]
After a finite ramification, this module has a full frame of rank-one
formal exponential functionals. If $a_D$ is the leading coefficient
of $\varphi$, every zero-exponential functional in this frame takes
the seed $[1]$ to a series of valuation at least
\begin{equation}\label{eq:valuationbound}
 v(h)-\frac{v(a_D)}D.
\end{equation}
In particular, if $Dv(h)>v(a_D)$, the differential submodule generated
by $[1]$ has no horizontal quotient taking $[1]$ to the constant one.
\end{theorem}
\begin{proof}
The two twisted derivatives commute, so the parameter connection
is well defined on the quotient. A nonzero polynomial of degree $m$
has spatial twisted derivative of degree $m+D-1$. Reduction of
successive leading terms therefore gives a basis
$1,x,\ldots,x^{D-2}$ and dimension $D-1$.

Split $\varphi_x$, retaining root multiplicities. For a critical root
$\xi$, expand
\[
 \varphi(\xi+u)=c_\xi+\sum_{j=2}^D a_{\xi,j}u^j,
 \qquad
 \rho_\xi=\max_{2\leq j\leq D,\ a_{\xi,j}\ne0}
                    \frac{-v(a_{\xi,j})}{j}.
\]
Enlarge the ramification so that all these scales are integral, and
set $\delta_\xi=r^{\rho_\xi}$. The polynomial
\[
 \varphi(\xi+\delta_\xi y)-c_\xi
\]
has integral coefficients, no linear term, and at least one unit
coefficient. Its reduction $F_\xi(y)$ has degree $d_\xi\geq2$.
Here and below we may replace $r$ by a root of $r$ and renormalize
the valuation; every displayed strict inequality is unchanged.
As $a_{\xi,D}=a_D$,
\begin{equation}\label{eq:balllower}
 \rho_\xi\geq-v(a_D)/D.
\end{equation}

\emph{The balls partition all critical roots.}
If another critical root $\zeta$ satisfies
$v(\zeta-\xi)\geq\rho_\xi$, translate the integral scaled polynomial
by $(\zeta-\xi)/\delta_\xi$. All its coefficients remain integral,
and the coefficient of degree $d_\xi$ remains a unit: higher-degree
terms have positive valuation. The linear term vanishes because
$\zeta$ is critical. It follows from the definition of the maximum
that $\rho_\zeta=\rho_\xi$. Thus their balls are identical.
If two of the balls intersect, one contains the other; the same
argument makes them identical. Otherwise they are disjoint.

The scaled derivative has reduction $F_\xi'$ of degree $d_\xi-1$.
It therefore has exactly $d_\xi-1$ integral roots, with multiplicity.
Here is a direct check of that count: factor it in a Puiseux field.
An integral root contributes a linear factor to the reduction;
for a root of negative valuation, factor out that root and the
remaining factor reduces to a constant. Multiplicativity of the
coefficient valuation leaves a nonzero constant scalar. Hence the
degree of the reduction counts exactly the integral roots.
The ball $B$ thus contains $d_B-1$ critical roots. Since
$\deg\varphi_x=D-1$,
\begin{equation}\label{eq:ballrankcount}
 \sum_B(d_B-1)=D-1.
\end{equation}

\emph{A functional for every dimension in a ball.}
The fixed polynomial $F_B$ has twisted quotient
\[
 H_B=\C[y]/(\partial_y+F_B')\C[y]
\]
with basis $1,y,\ldots,y^{d_B-2}$. Let $\ell^0_{B,j}$ be its dual
basis. Write
\[
 \varphi(\xi_B+\delta_B y)=c_B+F_B(y)+E_B(r,y),
 \qquad E_B\in r\C[[r]][y].
\]
For a polynomial amplitude $q$ define
\begin{equation}\label{eq:ballfunctional}
 \ell_{B,j}(q)=h\delta_B e^{c_B}
   \ell^0_{B,j}\!\left(e^{E_B(r,y)}q(\xi_B+\delta_B y)\right).
\end{equation}
Each coefficient uses finitely many polynomial terms. The exponential
in this formula is purely formal. The product rule conjugates the
spatial twisted derivative to $\partial_y+F_B'$, so the functional
annihilates the defining relations of $M$.

We spell out the parameter check because it retains the moving volume.
At fixed $y$, put $w=\xi_B'+\delta_B'y$. Differentiating
\eqref{eq:ballfunctional}, then subtracting its value on
$q_r+(\varphi_r+h'/h)q$, leaves, inside the fixed functional,
\[
 e^{E_B}\left(wq_x+\frac{\delta_B'}{\delta_B}q+\varphi_xwq\right)
 =\frac1{\delta_B}(\partial_y+F_B')
                      \left(e^{E_B}wq\right).
\]
It is zero. Thus the functionals are horizontal, including the
Jacobian $\delta_B$ and normalization $h$.

\emph{Completeness and the diagonal connection.}
Choose the polynomials $q_{B,j}$ of
Lemma~\ref{lem:ballseparation}, for $0\leq j<d_B-1$.
Index rows by functionals $(B,j)$ and columns by amplitudes $(C,k)$,
and form the normalized evaluation matrix
\[
 A_{(B,j),(C,k)}=
   \frac{\ell_{B,j}(q_{C,k})}{h\delta_Be^{c_B}}.
\]
Its entries are integral, and its reduction is the identity:
$E_B$ has positive order and the amplitudes reduce to the chosen
monomials on their own ball and zero on every other ball. Thus $A$
is invertible over the ramified Laurent field. There are $D-1$ rows
by \eqref{eq:ballrankcount}, equal to $\dim M$, so the amplitude
classes are a basis and these functionals account for the entire module.

Remove from $c_B$ its finite negative-power part $q_B(r)$.
The scalar
\[
 g_B=h\delta_B\exp(c_B-q_B)
\]
is an invertible Laurent series: $\exp(c_B-q_B)$ is a formal
power-series unit, with a nonzero constant term. Consequently the
matrix of the functionals $e^{-q_B}\ell_{B,j}$ in the amplitude
basis is $\operatorname{diag}(g_B)A$, with $g_B$ repeated for every
row on ball $B$. It is invertible over the same Laurent field.
The horizontality already proved says precisely that these rows
give an isomorphism to the direct sum of rank-one connections
$\partial_r+q_B'$, one for each pair $(B,j)$. This supplies the diagonal
gauge explicitly; it does not infer semisimplicity from a rank count.
If $q_B=0$, formula~\eqref{eq:ballfunctional} takes the seed to a
series of valuation at least $v(h)+\rho_B$, proving
\eqref{eq:valuationbound}.

Every rank-one differential module is simple, so this constructed
direct sum is semisimple. A quotient on the submodule generated by
the seed therefore extends to the whole module by choosing a
differential complement. A nonzero exponential factor has no
horizontal map to the Laurent field: the logarithmic derivative of
a nonzero Laurent series has at most a simple pole, whereas $q_B'$
has a pole of order at least two when $q_B\ne0$. On the zero part,
every horizontal functional is a constant combination of those in
the full frame. If the seed bound is strictly positive, their values
cannot sum to one.
\end{proof}

The argument is a formal stationary-phase construction for polynomial
phases over a Laurent field, including nonconvergent coefficients.
For the classical local Fourier--Laplace theory of exponential
connections, see Sabbah~\cite{SabbahStationary}.

\begin{corollary}[Leading-coefficient obstruction]\label{cor:leadingobstruction}
For a rational polynomial phase $\varphi(t,x)$ of degree $D\geq2$
and nonzero rational scalar normalization $h(t)$, let $v_\infty$
mean order at $\varepsilon=1/t=0$. If
\begin{equation}\label{eq:leadingcriterion}
 v_\infty(h^D/a_D)>0,
\end{equation}
its seed submodule cannot have a rational horizontal quotient with
seed value one. In particular any horizontal formal functional at
another point cannot take that seed to one.
\end{corollary}
\begin{proof}
Apply the theorem at infinity. A rational horizontal quotient remains
a quotient on extension to the formal Laurent field. A formal
functional with seed value one restricts to such a rational quotient:
its values on successive seed derivatives are $1,0,0,\ldots$.
\end{proof}

The inequality must be strict. For example
$\varphi=-x^2/(2t^2)$ and $h=1/t$ have
$v_\infty(h^2/a_2)=0$, and their normalized Gaussian period is a
nonzero constant. Rescaling that functional gives seed value one,
so the conclusion would fail if equality were allowed.


\paragraph{AI assistance and review.}
ChatGPT (OpenAI) assisted with mathematical arguments, literature checks,
exposition, proof review and verification code.

\begin{thebibliography}{99}
\small
\bibitem{Long}
C.~D.~Long, \emph{Small Counterexamples to the Gaussian Moments Conjecture},
public preprint (2026),
\href{https://arxiv.org/abs/2607.18186v1}{arXiv:2607.18186v1}.
\bibitem{Wilson}
M.~Wilson, \emph{A face-isolation proof of the two-variable Gaussian
Moments Conjecture}, public preprint (2026),
\href{https://arxiv.org/abs/2607.23887v1}{arXiv:2607.23887v1}.
\bibitem{Liu}
Z.~Liu, \emph{Sharp finite Gaussian moments for quadratic polynomials},
version~3, public preprint (2026),
\href{https://doi.org/10.5281/zenodo.21688251}{doi:10.5281/zenodo.21688251}.

\bibitem{SabbahStationary}
C.~Sabbah, \emph{An explicit stationary phase formula for the local
formal Fourier--Laplace transform}, in \emph{Singularities I},
Contemporary Mathematics \textbf{474}, American Mathematical Society,
2008, 309--330;
\href{https://arxiv.org/abs/0706.3570v2}{arXiv:0706.3570v2}.

\bibitem{MannaaCoquand}
B.~Mannaa and T.~Coquand, \emph{Dynamic Newton--Puiseux Theorem},
public preprint (2013), Corollary~1.10,
\href{https://arxiv.org/abs/1304.6770v2}{arXiv:1304.6770v2}.

\end{thebibliography}
\end{document}

